% Row retrieval poster, A0 portrait (workshop). Design page 1000 x 1414.2 mm, printed at 0.841 by poster_A0.tex.
% PROVENANCE: figures/fig_overview.pdf = the paper's Fig. 1; figures/fig_scaling.pdf = the paper's Fig. 2;
% tables/results.tex = the manuscript's table, unchanged. Prose pasted from paper/main.tex.
\documentclass{article}
\usepackage[paperwidth=1000mm, paperheight=1414.2mm, margin=0mm]{geometry}   % A0 aspect; printed at 0.841
\usepackage{fontspec}
\IfFontExistsTF{Helvetica Neue}{%
  \setmainfont{Helvetica Neue}[UprightFont=* Light, BoldFont=* Medium, ItalicFont=* Light Italic,
    BoldItalicFont=* Medium Italic]
  \newfontfamily\heavy{Helvetica Neue Bold}
  \newfontfamily\thin{Helvetica Neue Thin}
}{% fallback where Helvetica Neue is not installed (Suggested, not yet tested in real work)
  \IfFontExistsTF{texgyreheros-regular.otf}{\setmainfont{texgyreheros}[Extension=.otf, UprightFont=*-regular,
    BoldFont=*-bold, ItalicFont=*-italic, BoldItalicFont=*-bolditalic]}{\setmainfont{lmsans10}[Extension=.otf,
    UprightFont=*-regular, BoldFont=*-bold, ItalicFont=*-oblique, BoldItalicFont=*-boldoblique]}
  \newcommand{\heavy}{\bfseries}\newcommand{\thin}{}
}
\usepackage{amsmath}
\usepackage[mathrm=sym]{unicode-math}
% maths is set as in the paper: Computer Modern (Latin Modern Math, also for \mathrm), at the text's x-height
\setmathfont{latinmodern-math.otf}[Scale=MatchLowercase]
\usepackage{microtype}
\usepackage{graphicx}
\usepackage[table]{xcolor}
\usepackage{booktabs,array,tabularx}
\usepackage{tikz}
\usetikzlibrary{calc,positioning,fit}
\pagestyle{empty}
\setlength{\parindent}{0pt}

% ---------- poster chrome (figures keep the paper's own colors) ----------
\definecolor{ink}{HTML}{1A1A1A}
\definecolor{slate}{HTML}{4A4F5A}
\definecolor{accent}{HTML}{2B50AA}   % the one accent: the method's name, the coined term, key phrases. Use the
                                     % host's or the paper's color here.

% ---------- type scale (design pt; printed = x 0.841) ----------
\newcommand{\sectfont}{\fontsize{46}{54}\selectfont\bfseries}
\newcommand{\absfont}{\fontsize{33}{42.5}\selectfont}
\newcommand{\bodyfont}{\fontsize{27}{35}\selectfont}
\newcommand{\captfont}{\fontsize{25}{31}\selectfont}
\newcommand{\mathfont}{\fontsize{27}{34}\selectfont}

% ---------- a figure of the paper, or a labelled grey box while the file is missing ----------
% \fig[includegraphics options]{width cm}{placeholder height cm}{file}
\newcommand{\fig}[4][]{\IfFileExists{#4}{\includegraphics[width=#2cm,#1]{#4}}{%
  {\setlength{\fboxsep}{0pt}\setlength{\fboxrule}{0.06cm}\fcolorbox{slate}{black!4}{%
  \begin{minipage}[c][#3cm][c]{\dimexpr#2cm-0.12cm\relax}\centering\captfont\color{slate}%
  [figure: \detokenize{#4}]\end{minipage}}}}}

% ---------- the manuscript's table, unchanged, as a plain box ----------
% \resultstable{width}: for a table that sets itself to \linewidth (tabularx, or \resizebox{\linewidth}{!}{...})
% \resultstablescaled{width}: for a plain tabular at its natural width, scaled as a whole to the strip's width
\newcommand{\resultstable}[1]{\begingroup
  \renewenvironment{table}[1][]{}{}\renewcommand{\caption}[2][]{}\renewcommand{\label}[1]{}%
  \begin{minipage}{#1}\begin{table}[t]
\centering
\caption{\textbf{Row retrieval ranks first on average.} Accuracy on four dataset groups (mean over datasets $\pm$ SD over five seeds) and mean rank over all 48 datasets; lower rank is better. Bold marks the highest value in a column.}
\label{tab:results}
\small
\begin{tabular}{lccccc}
\toprule
method & small & medium & large & very large & mean rank $\downarrow$\\
\midrule
Gradient-boosted trees & .812{\scriptsize$\pm$.004} & .843{\scriptsize$\pm$.003} & .861{\scriptsize$\pm$.002} & \textbf{.874}{\scriptsize$\pm$.002} & 2.4\\
Base model, random context & \textbf{.826}{\scriptsize$\pm$.005} & .839{\scriptsize$\pm$.004} & .842{\scriptsize$\pm$.004} & .843{\scriptsize$\pm$.003} & 3.1\\
Base model, fine-tuned & .824{\scriptsize$\pm$.006} & .846{\scriptsize$\pm$.004} & .858{\scriptsize$\pm$.003} & .866{\scriptsize$\pm$.003} & 2.9\\
\midrule
Row retrieval (ours) & \textbf{.826}{\scriptsize$\pm$.005} & \textbf{.852}{\scriptsize$\pm$.003} & \textbf{.866}{\scriptsize$\pm$.002} & .871{\scriptsize$\pm$.002} & \textbf{1.6}\\
\bottomrule
\end{tabular}
\end{table}
\end{minipage}\endgroup}
\newsavebox{\tablebox}
\newcommand{\resultstablescaled}[1]{\begingroup
  \renewenvironment{table}[1][]{}{}\renewcommand{\caption}[2][]{}\renewcommand{\label}[1]{}%
  \renewcommand{\centering}{}\sbox{\tablebox}{\begin{table}[t]
\centering
\caption{\textbf{Row retrieval ranks first on average.} Accuracy on four dataset groups (mean over datasets $\pm$ SD over five seeds) and mean rank over all 48 datasets; lower rank is better. Bold marks the highest value in a column.}
\label{tab:results}
\small
\begin{tabular}{lccccc}
\toprule
method & small & medium & large & very large & mean rank $\downarrow$\\
\midrule
Gradient-boosted trees & .812{\scriptsize$\pm$.004} & .843{\scriptsize$\pm$.003} & .861{\scriptsize$\pm$.002} & \textbf{.874}{\scriptsize$\pm$.002} & 2.4\\
Base model, random context & \textbf{.826}{\scriptsize$\pm$.005} & .839{\scriptsize$\pm$.004} & .842{\scriptsize$\pm$.004} & .843{\scriptsize$\pm$.003} & 3.1\\
Base model, fine-tuned & .824{\scriptsize$\pm$.006} & .846{\scriptsize$\pm$.004} & .858{\scriptsize$\pm$.003} & .866{\scriptsize$\pm$.003} & 2.9\\
\midrule
Row retrieval (ours) & \textbf{.826}{\scriptsize$\pm$.005} & \textbf{.852}{\scriptsize$\pm$.003} & \textbf{.866}{\scriptsize$\pm$.002} & .871{\scriptsize$\pm$.002} & \textbf{1.6}\\
\bottomrule
\end{tabular}
\end{table}
}\resizebox{#1}{!}{\usebox{\tablebox}}\endgroup}

% ---------- grammar: section title + rule, then pieces chained below ----------
% Text is fully justified. TikZ sets \spaceskip from the 10 pt font before the font switch inside text-width
% nodes; \txt resets it so word spaces come from the poster font. \hyphenpenalty 1000: 5000 left loose lines
% and one-word last lines.
\tikzset{
  piece/.style={anchor=north west, inner sep=0pt, outer sep=0pt},
  sect/.style={piece, font=\sectfont, text=ink, text height=1.22cm, text depth=0.30cm},
}
\newcommand{\rulebelow}[2]{\fill[ink] ($(#1.south west)+(0,-0.3)$) rectangle ++(#2,-0.09);}
\newcommand{\txt}[3]{\begin{minipage}[t]{\linewidth}#1\spaceskip=0pt\xspaceskip=0pt\frenchspacing
  \hyphenpenalty=1000\exhyphenpenalty=1000\emergencystretch=2em\color{#2}#3\end{minipage}}
\newcommand{\body}[1]{\txt{\bodyfont}{ink}{#1}}
\newcommand{\capt}[1]{\txt{\captfont}{slate}{#1}}
\newcommand{\hi}[1]{\textbf{\color{accent}#1}}
\renewcommand{\cite}[1]{{\color{slate}[#1]}}
\newcommand{\bignum}[1]{{\thin\fontsize{72}{72}\selectfont\color{accent}#1}}
\newcommand{\numlab}[1]{{\captfont\color{slate}#1}}
% label under a brace: every label gets one height and depth, so all labels share a baseline
\newcommand{\lab}[1]{\text{\fontsize{22}{26}\selectfont\color{slate}\vphantom{fp}#1}}
\newcommand{\datatable}[1]{\begingroup\fontsize{27}{32}\selectfont\spaceskip=0pt\renewcommand{\arraystretch}{1.28}%
  \setlength{\heavyrulewidth}{2pt}\setlength{\lightrulewidth}{1.1pt}\setlength{\aboverulesep}{0.18em}%
  \setlength{\belowrulesep}{0.22em}#1\endgroup}

\begin{document}
\begin{tikzpicture}[remember picture, overlay, x=1cm, y=1cm]
\coordinate (O) at (current page.south west);
\newcommand{\at}[2]{($(O)+(#1,#2)$)}

% =====================================================================
% HEADER (black on white): title, authors, affiliations, a full-width rule
% =====================================================================
\node[piece] (title) at \at{3}{139.32} {\begin{minipage}[t]{72cm}\raggedright
  \fontsize{88}{98}\selectfont\heavy\textcolor{accent}{Row Retrieval} Makes Tabular Foundation\\Models Scale with Context\end{minipage}};
\node[piece, below=0.6cm of title.south west, anchor=north west] (authors) {\fontsize{35}{42}\selectfont
  \underline{Jane Doe}\textsuperscript{1}, John Roe\textsuperscript{1}};
\node[piece, below=0.45cm of authors.south west, anchor=north west, text width=94cm] (affil) {%
  \begin{minipage}[t]{94cm}\raggedright\fontsize{23}{29}\selectfont\spaceskip=0pt\xspaceskip=0pt\color{slate}%
  \mbox{\textsuperscript{1}Department of Computer Science, Example University}\end{minipage}};
\coordinate (hdr) at ($(affil.south west)+(0,-0.55)$);
\fill[ink] (hdr) rectangle ($(hdr)+(94,-0.2)$);

% =====================================================================
% BAND 1: Abstract (55.5 cm) | Key idea: the paper's concept figure and its defining equations (36.5 cm)
% =====================================================================
\def\Aw{55.5}
\node[sect, below=1.0cm of hdr, anchor=north west] (s1) {Abstract};
\rulebelow{s1}{\Aw}
\node[piece, below=0.75cm of s1.south west, anchor=north west, text width=\Aw cm] (abs) {\txt{\absfont}{ink}{%
  \textbf{Can a tabular foundation model use more training rows than fit in its context?} In-context learners for
  tables condition on a fixed number of labelled rows, and their accuracy stops improving once that budget is
  reached. We introduce \hi{row retrieval}, which selects for every query row the training rows \hi{most similar to
  it} in the model's own embedding space and places only those in the context. Retrieval makes in-context tabular
  models scale far beyond their context window at almost no cost. Row retrieval requires \hi{no fine-tuning} and
  applies to any in-context tabular model.}};

\def\Bx{60.5}\def\Bw{36.5}
\node[sect] (s2) at ($(s1.north west)+(\Bx-3,0)$) {Row retrieval};
\rulebelow{s2}{\Bw}
\node[piece, below=0.75cm of s2.south west, anchor=north west, text width=\Bw cm] (defs) {%
  \begin{minipage}[t]{\Bw cm}\raggedright\mathfont\spaceskip=0pt\color{ink}%
  {\captfont\color{slate}the context of a query row $x$}\\[0.3em]
  {\setlength{\fboxrule}{2.2pt}\setlength{\fboxsep}{0.27cm}\fcolorbox{accent}{accent!7}{%
    \begin{minipage}{\dimexpr\Bw cm-2\fboxrule-2\fboxsep\relax}\centering
    {\fontsize{27}{32}\selectfont\bfseries\color{accent}row retrieval}%
    \fontsize{28}{34}\selectfont\abovedisplayskip=0.3em\belowdisplayskip=0.4em
    \[\mathcal{C}(x)=\operatorname{top-}k_{\,i}\ \frac{\langle f_\theta(x), f_\theta(x_i)\rangle}{\lVert f_\theta(x)\rVert\,\lVert f_\theta(x_i)\rVert}\]
    {\captfont\color{slate}The index is built once per dataset and reused for every query.}\end{minipage}}}\end{minipage}};

% =====================================================================
% BAND 2: Method, the paper's two method figures side by side (2 x 45.85 cm, gutter 2.3 cm)
% =====================================================================
\node[fit=(abs)(defs), inner sep=0pt] (b1) {};
\node[sect] (s3) at ($(b1.south west)+(0,-0.86)$) {Method};
\rulebelow{s3}{94}
\def\Fw{45.85}
\node[piece, below=0.75cm of s3.south west, anchor=north west] (f1) {\fig{\Fw}{16}{figures/fig_overview.pdf}};
\node[piece] (f2) at ($(f1.north west)+(48.15,0)$) {\fig[height=21.8cm, width=!]{\Fw}{16}{figures/fig_scaling.pdf}};
\node[fit=(f1)(f2), inner sep=0pt] (f12) {};
\node[piece, text width=\Fw cm] (c3) at ($(f1.south west |- f12.south)+(0,-0.6)$) {\capt{%
  \textbf{\color{ink}Row retrieval chooses the context per query.} (a) The base model fills its context with a
  random subsample of the training set. (b) Row retrieval embeds all training rows once, indexes them and fills the
  context with the nearest rows to each query.}};
\node[piece, text width=\Fw cm] (c4) at ($(f2.south west |- f12.south)+(0,-0.6)$) {\capt{%
  \textbf{\color{ink}Accuracy keeps rising with training rows only under retrieval.} Mean accuracy over the 48
  datasets against the number of available training rows. The dashed line marks the base model's context length.}};

% =====================================================================
% BAND 4: Data (31 cm): big numbers and the inventory | Results (61 cm): the manuscript's table and its key
% =====================================================================
\node[fit=(c3)(c4), inner sep=0pt] (b3) {};
\def\Lw{31}\pgfmathsetmacro\Rx{\Lw+5}\pgfmathsetmacro\Rw{92-\Lw}
\node[sect] (s6) at ($(b3.south west)+(0,-0.86)$) {Data};
\rulebelow{s6}{\Lw}
\node[piece, below=0.75cm of s6.south west, anchor=north west] (n6) {%
  \begin{tabular*}{\Lw cm}{@{\extracolsep{\fill}}ll@{}}
  \bignum{48} & \bignum{2,000,000}\\[0.3cm]
  \numlab{classification datasets} & \numlab{rows in the largest}
  \end{tabular*}};
\node[piece, below=0.7cm of n6.south west, anchor=north west, text width=\Lw cm] (t6) {\body{%
  On 48 classification datasets with 10,000 to 2,000,000 rows.}};

\node[sect, below=1.3cm of t6.south west, anchor=north west] (s7) {Evaluation};
\rulebelow{s7}{\Lw}
\node[piece, below=0.75cm of s7.south west, anchor=north west, text width=\Lw cm] (c7) {\body{%
  The gain grows with the size of the training set. On datasets with fewer than 8,192 rows, where the whole
  training set fits in the context, the two models are level.}};

\node[sect] (s9) at ($(s6.north west)+(\Rx-3,0)$) {Results};
\node[piece, below=0.23cm of s9.south west, anchor=north west] (t9) {\resultstablescaled{\Rw cm}};
\node[piece, below=-1.2cm of t9.south west, anchor=north west, text width=\Rw cm] (c9) {\body{%
  Row retrieval \textbf{ranks first on average} across the 48 datasets.}};
\node[piece, below=0.35cm of c9.south west, anchor=north west, text width=\Rw cm] (c9b) {\capt{%
  Accuracy on four dataset groups (mean over datasets $\pm$ SD over five seeds) and mean rank over all 48
  datasets; lower rank is better. Bold marks the highest value in a column.}};

% =====================================================================
% FOOTER (black on white): a full-width rule, references and contacts on the left, the release line on the right
% =====================================================================
\newcommand{\rsep}{\hspace{0.4em plus 0.2em minus 0.05em}}   % between references
\def\Fw{77.4}
\node[anchor=south east, inner sep=0pt] (npt) at \at{97}{1.3} {\fontsize{26}{31}\selectfont\bfseries
  Preprint coming soon};
\node[anchor=south, inner sep=0pt] (np) at ($(npt.north)+(0,0.35)$) {\fig{2.4}{2.4}{assets/no_photos.pdf}};
\node[anchor=south west, inner sep=0pt] (contacts) at \at{3}{1.3} {\begin{minipage}{\Fw cm}%
  \fontsize{22}{26}\selectfont\spaceskip=0pt\color{slate}%
  jane.doe@example.org\hfill Example University, Department of Computer Science\end{minipage}};
\node[anchor=south west, inner sep=0pt] (refs) at ($(contacts.north west)+(0,0.4)$) {\begin{minipage}{\Fw cm}%
  \fontsize{12}{15}\selectfont\spaceskip=0pt\frenchspacing\color{slate}\textbf{\color{ink}References}\rsep
  [1]~Chen and Guestrin, XGBoost, KDD~2016\rsep
  [2]~Malkov and Yashunin, HNSW, IEEE TPAMI~2020\end{minipage}};
\node[fit=(refs)(contacts)(np)(npt), inner sep=0pt] (foot) {};
\fill[ink] ($(O |- foot.north)+(3,0.55)$) rectangle ($(O |- foot.north)+(97,0.75)$);
\node[inner sep=0pt, minimum height=0.2cm, minimum width=94cm, anchor=south west] (footrule) at ($(O |- foot.north)+(3,0.55)$) {};

% ---------- log every box in page cm for box_check.py (add every new node here) ----------
\foreach \nm in {title,authors,affil,s1,abs,s2,defs,s3,f1,f2,c3,c4,s6,n6,t6,s7,c7,
                 s9,t9,c9,c9b,np,npt,refs,contacts,footrule}{
  \path let \p0=(O), \p1=(\nm.south west), \p2=(\nm.north east) in
    \pgfextra{\pgfmathsetmacro\xa{(\x1-\x0)/28.45276}\pgfmathsetmacro\ya{(\y1-\y0)/28.45276}%
      \pgfmathsetmacro\xb{(\x2-\x0)/28.45276}\pgfmathsetmacro\yb{(\y2-\y0)/28.45276}%
      \typeout{BOX \nm\space\xa\space\ya\space\xb\space\yb}};}
\end{tikzpicture}
\end{document}
